\documentclass[12pt]{article}
\usepackage[a4paper, margin=1in]{geometry}
\usepackage{graphicx}
\usepackage{caption}
\usepackage{parskip}
\title{Getting Started with the MATLAB Editor}
\author{}
\date{}

\begin{document}

\maketitle

\section*{1. What is MATLAB?}

MATLAB is a programming environment widely used for numerical computation, data analysis, and engineering tasks. It has an integrated editor where you can write, run, and save scripts (files containing code).

\section*{2. Exploring the MATLAB Interface}

When you open MATLAB, you’ll see several important windows:

\begin{itemize}
    \item \textbf{Command Window}: where you can type and immediately execute commands.
    \item \textbf{Editor}: where you write and save longer programs or scripts.
    \item \textbf{Current Folder}: shows your working directory (where files are saved).
\end{itemize}

\begin{figure}[h!]
    \centering
    \includegraphics[width=0.85\textwidth]{screenshots/matlab_windows.png}
    \caption{MATLAB interface: Command Window, Editor, Current Folder}
\end{figure}

\section*{3. Creating a New Script}

To start writing a program, click \texttt{New Script} in the toolbar. This opens the Editor window.

In the editor, type the following instruction:

\begin{verbatim}
disp('Hello, MATLAB!')
\end{verbatim}

This command displays a message in the Command Window when the script is run.

\begin{figure}[h!]
    \centering
    \includegraphics[width=0.85\textwidth]{screenshots/new_script.png}
    \caption{Writing a simple instruction in the Editor}
\end{figure}

\section*{4. Running the Script}

Click the \texttt{Run} button (green triangle) in the toolbar.

\textbf{What happens next?}
\begin{itemize}
    \item MATLAB will ask you to save the script first if it hasn’t been saved.
    \item Once saved, the code runs and the result is shown in the \textbf{Command Window}.
\end{itemize}

\begin{figure}[h!]
    \centering
    \includegraphics[width=0.85\textwidth]{screenshots/run_script.png}
    \caption{Running the script — output appears in the Command Window}
\end{figure}

\section*{5. Saving the Script}

Click \texttt{File > Save} or the disk icon to save your script.

\begin{itemize}
    \item Choose a meaningful name, e.g. \texttt{hello.m}
    \item The \texttt{.m} extension is required for MATLAB scripts
\end{itemize}

\begin{figure}[h!]
    \centering
    \includegraphics[width=0.85\textwidth]{screenshots/save_script.png}
    \caption{Saving the script with a \texttt{.m} extension}
\end{figure}

\section*{6. Reopening a Saved Script}

To open a saved script again:

\begin{itemize}
    \item Find it in the \texttt{Current Folder} window and double-click it, or
    \item Use \texttt{File > Open} and select the file

\end{document}